\documentclass[a4paper,11pt]{article}
\usepackage[utf8]{inputenc}
\usepackage[T1]{fontenc}
\usepackage[french]{babel}
\usepackage[left=1cm,right=1cm,top=.5cm,bottom=0cm]{geometry}
\usepackage{enumitem}
\usepackage{hyperref}
\usepackage{xcolor}
\usepackage{titlesec}
\usepackage{fontawesome5}
\usepackage{lmodern}
\usepackage[normalem]{ulem}

% --- Couleurs Eiffage Rail ---
\definecolor{primary}{RGB}{30, 90, 50}
\definecolor{secondary}{RGB}{80, 80, 80}

% --- Config Liens ---
\hypersetup{colorlinks=true, linkcolor=primary, filecolor=primary, urlcolor=primary}

% --- Titres ---
\titleformat{\section}{\large\bfseries\color{primary}\uppercase}{}{0em}{}[\titlerule]
\titlespacing{\section}{0pt}{8pt}{5pt}

% --- Commandes ---
\newcommand{\resumeItem}[1]{\item\small{#1}}
% NOTE: NE PAS utiliser resumeSubheading avec tabular* car cela cause des superpositions de texte
% Utiliser le format avec \hfill et \\ pour les retours à la ligne

\begin{document}

% --- EN-TÊTE ---
\begin{center}
    {\Large \textbf{Loïc Aron MBASSI EWOLO}} \\ \vspace{4pt}
    \textbf{\large Alternance / Stage PFE -- Développement Logiciel Applicatif} \\
    \textit{\small Disponible dès avril 2026 (stage 4-6 mois) ou septembre 2026 (alternance)} \\ \vspace{4pt}
    \small
    \faMapMarker\ Cergy, France (Mobile IDF -- Élancourt) \quad
    \faPhone\ (+33) 7 59 52 33 98 \quad
    \faEnvelope\ \href{mailto:loicaron.mbassiewolo@ensea.fr}{loicaron.mbassiewolo@ensea.fr} \\ \vspace{2pt}
    \faLinkedin\ \href{https://www.linkedin.com/in/loic-mbassi/}{linkedin.com/in/loic-mbassi} \quad
    \faGithub\ \href{https://github.com/Nameless0l}{github.com/Nameless0l} \quad
    \faGlobe\ \href{https://mbassiloic.tech}{mbassiloic.tech}
\end{center}

\vspace{-6pt}

% --- PROFIL ---
\noindent\small{Deux ans de maintenance et d'évolution d'une plateforme de télémédecine complète (SOSDOCTA, Laravel, MySQL, RGPD) et une expérience en développement backend fintech (CSC, Docker, Redis) constituent une base directe pour les missions de maintenance corrective et évolutive décrites. Maîtrise de Java, Python, SQL/NoSQL et des architectures microservices (GoFind : MongoDB, RabbitMQ, Docker). Disponible dès avril 2026 pour Élancourt.}

% --- FORMATION ---
\section{Formation}
\begin{itemize}[leftmargin=*, label=\tiny$\bullet$, nosep]
    \item \textbf{ENSEA -- École Nationale Supérieure de l'Électronique et de ses Applications} \hfill \textit{2025 -- 2027} \\
    \small{Double diplôme d'Ingénieur -- Systèmes Embarqués, Signal, Intelligence Artificielle.}
    \item \textbf{École Polytechnique de Yaoundé (1\textsuperscript{ère} école d'ingénieur CEMAC)} \hfill \textit{2021 -- 2025} \\
    \small{Diplôme d'Ingénieur de Conception (Master 2) : \textbf{Mathématiques Appliquées}, Génie Logiciel, Algorithmique.}
\end{itemize}

% --- COMPÉTENCES ---
\section{Compétences Techniques}
\begin{itemize}[leftmargin=*, nosep]
    \item \small{\textbf{Langages :} Java, Python, PHP/Laravel, JavaScript, C/C++}
    \item \small{\textbf{Bases de données :} MySQL, PostgreSQL (SQL) -- MongoDB, Redis (NoSQL)}
    \item \small{\textbf{Frameworks \& Outils :} Laravel, NestJS, Spring Boot, Docker, Git/GitHub, REST APIs}
    \item \small{\textbf{Méthodes :} Agile/Scrum, CI/CD, Microservices, Documentation technique}
    \item \small{\textbf{Langues :} Français (Natif), Anglais (Professionnel -- Technique)}
\end{itemize}

% --- EXPÉRIENCES / PROJETS ---
\section{Expériences \& Projets}
\begin{itemize}[leftmargin=*, label={-}]

\item \textbf{Développeur Web Full Stack -- SOSDOCTA} \hfill \textit{Fév 2022 -- Mars 2024} \\
\textit{Remote (Belgique / Canada) -- Plateforme de Télémédecine} \hfill \textit{Laravel, PHP, MySQL, OAuth2/JWT}
\vspace{-2pt}
    \begin{itemize}[leftmargin=0.4cm, nosep, label=\tiny$\bullet$]
        \item \small{Développement et maintenance sur 2 ans d'une plateforme complète (médecins, patients, réservations, paiements) -- directement comparable aux missions de maintenance corrective et évolutive d'Eiffage Rail.}
        \item \small{Conception et évolution de dashboards de suivi en temps réel (traçabilité des rendez-vous, alertes, historique patients) -- même logique que la gestion de matériel et d'alertes de maintenance.}
        \item \small{Gestion et structuration de bases de données MySQL, conformité RGPD, sécurisation des données sensibles.}
    \end{itemize}
\vspace{3pt}

\item \textbf{Développeur Backend -- CSC (Cameroon Software Company)} \hfill \textit{Oct 2024 -- Jan 2025} \\
\textit{Yaoundé -- Fintech / Bancaire} \hfill \textit{Laravel, MySQL, RESTful APIs, Docker, Redis, Git}
\vspace{-2pt}
    \begin{itemize}[leftmargin=0.4cm, nosep, label=\tiny$\bullet$]
        \item \small{Conception backend d'une plateforme de transactions sécurisée : gestion des accès, concurrence, conformité normes financières.}
        \item \small{Mise en place d'APIs RESTful robustes et d'une architecture scalable. Utilisation de Redis pour la gestion des sessions et du cache. Déploiement Docker.}
    \end{itemize}
\vspace{3pt}

\item \textbf{GoFind -- Architecture Microservices} \hfill \textit{2024} \\
\textit{Projet Académique} \hfill \textit{Laravel, NestJS, MongoDB, RabbitMQ, Docker, API Gateway}
\vspace{-2pt}
    \begin{itemize}[leftmargin=0.4cm, nosep, label=\tiny$\bullet$]
        \item \small{Conception d'une architecture microservices complète avec communication asynchrone via RabbitMQ et API Gateway centralisée.}
        \item \small{Utilisation de MongoDB (NoSQL) pour la gestion de données non structurées et de Laravel pour les services métier. Orchestration Docker.}
    \end{itemize}
\vspace{3pt}

\item \textbf{Application de Gestion -- Projet Java (Design Patterns)} \hfill \textit{2023} \\
\textit{Projet Académique (POO)} \hfill \textit{Java, Swing/JavaFX, Design Patterns (Singleton, Factory, Observer)}
\vspace{-2pt}
    \begin{itemize}[leftmargin=0.4cm, nosep, label=\tiny$\bullet$]
        \item \small{Développement d'une application desktop robuste en Java avec interface graphique Swing/JavaFX. Application intensive des patterns de conception (Singleton, Factory, Observer) pour la maintenabilité du code.}
        \item \small{Architecture orientée objet rigoureuse, adaptée aux contextes de développement et de maintenance d'applications internes en entreprise.}
    \end{itemize}
\vspace{3pt}

\end{itemize}

% --- DISTINCTIONS ---
\section{Distinctions \& Leadership}
\begin{itemize}[leftmargin=*, nosep, label=\tiny$\bullet$]
    \item \small{\textbf{1\textsuperscript{ère} Place} -- CITS National Hackathon 2024 (Smart City, 75 équipes) -- optimisation de collecte et d'exploitation de données terrain.}
    \item \small{\textbf{1\textsuperscript{ère} Place} -- IndabaX Africa 2025 (IA Santé) : déploiement d'API, traitement de données.}
    \item \small{\textbf{Head of Projects Cell \& Lead AI Cell} (ENSPY) : coordination d'équipes, animation workshops, reporting régulier. +50\% membres.}
    \item \small{\textbf{Laravel API Generator} (Open Source) : outil de génération de code automatique, +30 étoiles GitHub, déployé sur Packagist.}
\end{itemize}

\end{document}
